Qwen2-VL、Qwen2-Audio、Pixtral系、Ai2のmultimodal系などは、vision/audioを扱うmodel familyが急速に増えたことを示す。ただし入力modalitiesを理解できることと、音声をnativeに生成しながら会話できること、さらに画像・動画を生成することは別の能力である。2024-H2はその境界が明確になり始めた。\autocite{src-56b3abdc51a9,src-72d4a57e9c16,src-9f35f00f9f9c,src-d483c7d9838b,src-feb76866f5df}


OpenAIのRealtime APIとKyutaiのMoshiは、speechを単なる入出力変換ではなくinteraction loopとして扱う方向を示した。一方はproduct/API surface、もう一方は研究・公開systemとしての性格が異なるため、同じ『音声AI』としてavailabilityやdeploymentを統合しない。full-duplex、latency、native audioといった論点は、model architectureとtransport/product surfaceの両方にまたがる。\autocite{src-4a870e1e2c95,src-c5ed64e39d0c}


FLUX、Stable Diffusion 3.5、CogVideoX、Hunyuan系、LTX-Video、Movie Gen、SoraやVeo系の動きは、生成mediaの競争がimageからvideoへ大きく広がったことを示す。ただしresearch announcement、open-weight ecosystem、hosted product rolloutは同じavailability classではない。特にopen videoは、weightsの公開だけでなくVRAM、runtime、sampling costなどlocal deployment条件を含めて評価する必要がある。\autocite{src-06d2114e2823,src-2d134659d4de,src-d1b80acf98c2,src-d33cd3d41599,src-d50c19922f51,src-e91cbf23fde6,src-f6c8da8253fb}


同時期の複数研究は、画像・音声・動画を別々のspecialized pipelineで扱うだけでなく、tokenization、diffusion/flow、joint representationなどを通じて統合的に扱う設計空間を探索した。これらは将来の統合multimodal systemを示唆するが、paper上のauthor-reported evaluationをそのままproduction availabilityや優位性へ変換しない。\autocite{src-25c4073c2a29,src-5ff9bff435a6,src-680a61ed20c0,src-770e33132b09,src-e38a55668df9,src-e7d2578ccb4d,src-f6b059ed7fb2,src-fa0d8178b72a}


